\section{Conclusions}
\label{sec:conclusions}

We have demonstrated through systematic three-dimensional random walk simulations that the viscoelastic gel-point of a percolating network is not a fixed property of occupation density — it is determined by the growth dynamics of the forming cluster, and can be continuously programmed via a single nucleation-density parameter.

Three main results are established. First, for random (Bernoulli) percolation networks, the Winter--Chambon gel criterion ($\alpha = 0.5$) is satisfied at $\pcprime \approx 1 - \pc \approx 0.683$, quantitatively identifying the rheological gel-point with the loss of spanning connectivity in the void phase. This is the natural gel-point baseline for any network-forming material that grows without spatial correlation.

Second, replacing random growth with adjacency-constrained (templated) cluster growth shifts the gel-point by $\Delta\pcprime = +0.20$ to $\pcprime \approx 0.886$, and changes the character of the transition: the critical region is 3.4-times broader and probe particles remain persistently sub-diffusive above the gel-point — a morphological signature of the more open, connected void channels maintained by compact cluster growth.

Third, a nucleation-density parameter $\kappa$ provides continuous, monotonic interpolation between the Bernoulli and templated limits, spanning a total tunable range $\Delta\pcprime \approx 0.31$ ($\pcprime \in [0.683, 0.993]$). The $\kappa$ parameter maps directly onto the fraction of growing sites that originate as independent nuclei, providing a physically motivated handle that may be realised experimentally through nucleation site density, cross-linker concentration, or thermal ramp rate.

Frequency-domain viscoelastic spectra derived via the GSER confirm the $\delta(\omega) = 45°$ Winter--Chambon signature at $\pcprime$ for both gelation classes, establishing a direct bridge from simulation to measurable microrheology.

Together these results provide a conceptual and quantitative framework for programming viscoelastic gel-points in network-forming materials through control of cluster growth dynamics rather than chemical composition alone.

